%%%%%%%%%%%%%%%%%
% This is an sample CV template created using altacv.cls
% (v1.7.4, 30 July 2025) written by LianTze Lim (liantze@gmail.com). Compiles with pdfLaTeX, XeLaTeX and LuaLaTeX.
%
%% It may be distributed and/or modified under the
%% conditions of the LaTeX Project Public License, either version 1.3
%% of this license or (at your option) any later version.
%% The latest version of this license is in
%%    http://www.latex-project.org/lppl.txt
%% and version 1.3 or later is part of all distributions of LaTeX
%% version 2003/12/01 or later.
%%%%%%%%%%%%%%%%

\documentclass[10pt,a4paper,withhyper]{altacv}

% Layout and typography
\geometry{left=1.35cm,right=1.35cm,top=1.35cm,bottom=1.35cm}

\iftutex
  \setmainfont{Roboto Slab}
  \setsansfont{Lato}
  \renewcommand{\familydefault}{\sfdefault}
\else
  \usepackage[rm]{roboto}
  \usepackage[defaultsans]{lato}
  \renewcommand{\familydefault}{\sfdefault}
\fi
\usepackage{microtype}
\DisableLigatures{encoding = *, family = *}

% Restrained aerospace-inspired palette
\definecolor{ResumeText}{HTML}{111827}
\definecolor{ResumeMuted}{HTML}{4B5563}
\definecolor{ResumeAccent}{HTML}{163A5F}
\definecolor{ResumeRule}{HTML}{CBD5E1}
\colorlet{name}{ResumeText}
\colorlet{tagline}{ResumeAccent}
\colorlet{heading}{ResumeAccent}
\colorlet{headingrule}{ResumeRule}
\colorlet{subheading}{ResumeText}
\colorlet{accent}{ResumeAccent}
\colorlet{emphasis}{ResumeText}
\colorlet{body}{ResumeMuted}

\renewcommand{\namefont}{\Huge\bfseries}
\renewcommand{\personalinfofont}{\footnotesize}
\renewcommand{\cvsectionfont}{\large\bfseries}
\renewcommand{\cvsubsectionfont}{\normalsize\bfseries}
\renewcommand{\cvItemMarker}{{\small\textbullet}}

\setlength{\parindent}{0pt}
\setlength{\parskip}{0pt}
\AtBeginEnvironment{itemize}{\small}

% Local helpers for a one-column ATS-friendly layout
\newcommand{\contactsep}{\enspace\textbar\enspace}
\newcommand{\resumesection}[1]{%
  \cvsection{#1}\vspace{-0.45\baselineskip}%
}
\newcommand{\resumeentry}[4]{%
  \textbf{#1}%
  \ifstrempty{#3}{}{ {\small\color{emphasis}(#3)}}%
  \par
  \ifstrempty{#2}{}{%
    {\emph{#2}}%
    \ifstrempty{#4}{}{%
      \hfill{\small\color{emphasis}#4}%
    }%
    \par
  }%
}
\newcommand{\skillgroup}[2]{%
  \textbf{#1}: #2\par
}

\begin{document}

% TODO: Confirm whether the candidate wants to keep the full Lattes URL in the header or move it to a shorter link label.
% TODO: Confirm whether the MBA should remain if the resume is later tightened further for one-page submissions.
% TODO: Confirm whether graph attention should remain attributed to the Loyal Wingmen project in all target applications, or only when the 2026 journal article is cited.

{\namefont\color{name} Davi Guanabara Arag\~ao\par}
{\large\bfseries\color{tagline} Engenheiro de Intelig\^encia Artificial e Aprendizado de M\'aquina (AI/ML) \textbar{} Pesquisador de Doutorado em Sistemas Aut\^onomos e Ci\^encia de Dados\par}
\vspace{0.35em}
{\small
S\~ao Paulo, Brasil\contactsep
\href{mailto:davi\_guanabara@live.com}{davi\_guanabara@live.com}\contactsep
+55 85 9 8673-6026\contactsep
\href{https://github.com/DaviGuanabara}{github.com/DaviGuanabara}\contactsep
\href{https://linkedin.com/in/davi-guanabara}{linkedin.com/in/davi-guanabara}\contactsep
\href{http://lattes.cnpq.br/2180140366316472}{lattes.cnpq.br/2180140366316472}\par
}
\vspace{0.55em}

\resumesection{Resumo Profissional}

Engenheiro de Intelig\^encia Artificial e Aprendizado de M\'aquina (AI/ML) e pesquisador de doutorado no Instituto Tecnol\'ogico de Aeron\'autica (ITA), atuando na aplica\c{c}\~ao de intelig\^encia artificial, aprendizado de m\'aquina e engenharia de software a sistemas aut\^onomos, computa\c{c}\~ao geoespacial e aplica\c{c}\~oes aeroespaciais. Atualmente contribui para iniciativas de pesquisa em colabora\c{c}\~ao com o SAC/ICEA, desenvolvendo m\'etodos computacionais para planejamento de espa\c{c}o a\'ereo urbano, tomada de decis\~ao inteligente e experimenta\c{c}\~ao baseada em simula\c{c}\~ao. Seus interesses de pesquisa concentram-se em aprendizado por refor\c{c}o, sistemas multiagentes, IA geoespacial e desenvolvimento de software cient\'ifico reprodut\'ivel para problemas complexos de engenharia.

\medskip
\resumesection{\'Areas de Especializa\c{c}\~ao}

\skillgroup{Sistemas Aut\^onomos}{
Tomada de Decis\~ao Aut\^onoma • Sistemas Multiagentes • Sistemas Aut\^onomos com UAVs
}

\skillgroup{Aprendizado de M\'aquina}{
Deep Learning • Aprendizado por Refor\c{c}o • Modelos Baseados em Aten\c{c}\~ao
}

\skillgroup{Engenharia de Software}{
Sistemas Reativos • Plataformas Experimentais
}

\skillgroup{Computa\c{c}\~ao Geoespacial}{
Modelagem Geoespacial • Modelagem Espa\c{c}o-Temporal
}

\medskip
\resumesection{Projetos Selecionados}

\resumeentry{Deep Reinforcement Learning with Graph Attention for Cooperative Loyal Wingmen}
{Programa de Pesquisa CTEDS | Brasil-Suécia ADS (Air Domain Study)}
{2024--2026}{}

\begin{itemize} 

\item Desenvolveu um framework de deep reinforcement learning para tomada de decis\~ao cooperativa de engajamento de amea\c{c}as em Loyal Wingmen UAVs operando sob observa\c{c}\~oes espa\c{c}o-temporais compartilhadas.

\item Prop\^os uma arquitetura neural com graph attention que permite compartilhamento adaptativo de informa\c{c}\~oes e aprendizado coordenado de pol\'iticas sob observabilidade parcial.

\item Projetou e avaliou experimentos baseados em simula\c{c}\~ao integrando Graph Neural Networks com multi-agent reinforcement learning para sistemas a\'ereos aut\^onomos.

\item Desenvolvido no \^ambito do Cooperative Threat Engagement with Heterogeneous Drone Swarms (CTEDS) research program, uma colabora\c{c}\~ao internacional Brasil--Su\'ecia envolvendo governo, academia e ind\'ustria.

\item Resultou em duas publica\c{c}\~oes com revis\~ao por pares:

      \emph{Journal of Intelligent \& Robotic Systems} (Springer Nature, 2026, DOI:\href{https://doi.org/10.1007/s10846-026-02415-8}{10.1007/s10846-026-02415-8}) e \emph{2024 Brazilian Symposium on Robotics (SBR) and Workshop on Robotics in Education}(DOI:\href{https://doi.org/10.1109/SBR/WRE63066.2024.10837749}{10.1109/SBR/WRE63066.2024.10837749}).
\end{itemize}

\resumeentry{SkyWeaver — Plataforma para Configuração Dinâmica do Espaço Aéreo para o BRUTM}
{Projeto de pesquisa em colabora\c{c}\~ao com SAC/ICEA}
{2025--Present}{}

\begin{itemize}
\item Projetou uma arquitetura reativa para Dynamic Airspace Configuration com o objetivo de contribuir para o conceito brasileiro de UTM (BRUTM).
\item Desenvolveu componentes modulares de software para gera\c{c}\~ao de cen\'arios geoespaciais, gerenciamento de estado operacional, roteamento e modelagem de infraestrutura de espa\c{c}o a\'ereo.
\item Investiga m\'etodos de otimiza\c{c}\~ao baseados em grafos para posicionamento de droneports, planejamento de rotas e provis\~ao de infraestrutura sob restri\c{c}\~oes operacionais.
\item A pesquisa em andamento explora t\'ecnicas de engenharia de software, otimiza\c{c}\~ao e IA para futuras opera\c{c}\~oes de UTM, U-Space e Advanced Air Mobility (AAM).
\end{itemize}

\resumeentry{OndeVale — Plataforma Inteligente para Avaliação de Imóveis}{Sistema de IA Implantado em Nuvem}{2025}{}
\begin{itemize}
\item Desenvolveu uma plataforma implantada em nuvem para avalia\c{c}\~ao de im\'oveis residenciais, combinando an\'alise geoespacial, machine learning e visualiza\c{c}\~ao web interativa.
\item Projetou o pipeline completo de machine learning, incluindo pr\'e-processamento de dados, feature engineering, treinamento de modelos, servi\c{c}os de infer\^encia e implanta\c{c}\~ao.
\item Construiu um backend em FastAPI e um frontend em Vue.js integrando visualiza\c{c}\~ao GIS com infer\^encia de modelo em tempo real.
\item Implantou a aplica\c{c}\~ao como um servi\c{c}o em nuvem acess\'ivel publicamente, conectando pr\'e-processamento de dados, infer\^encia de modelos, uma REST API e uma interface geoespacial interativa.
\end{itemize}

\resumeentry{Framework Reprodutível para Aprendizado Supervisionado:
Estudo de Caso em Previsão de Umidade Foliar}{}{2024--2025}{}
\begin{itemize}
\item Desenvolveu um modelo de deep learning para previs\~ao de umidade foliar em vinhedos com base em observa\c{c}\~oes meteorol\'ogicas do Iowa Environmental Mesonet.
\item Projetou uma arquitetura Deep Sets permutation-invariant com masked attention para modelar vizinhan\c{c}as espa\c{c}o-temporais heterog\^eneas.
\item Construiu um framework reprodut\'ivel de aprendizado supervisionado integrando stochastic contextual sampling, Bayesian hyperparameter optimization, bootstrap evaluation e gerenciamento determin\'istico de seeds.
\end{itemize}

\medskip
\resumesection{Forma\c{c}\~ao Acad\^emica}
\resumeentry{Doutorado em Engenharia Eletrônica e Computação (Informática)}{Instituto Tecnol\'ogico de Aeron\'autica (ITA)}{2025--Present}{}
\resumeentry{Especialização em Ciência de Dados (CEDS)}{Instituto Tecnol\'ogico de Aeron\'autica (ITA)}{2024--2026}{}
\resumeentry{Mestrado em Computação Aeronáutica (MPCOMP)}{Instituto Tecnol\'ogico de Aeron\'autica (ITA)}{2021--2024}{}
\resumeentry{Bacharelado em Engenharia de Controle e Automação}{Universidade de Fortaleza (UNIFOR)}{2013--2019}{}



\end{document}
